% --- Archon physics formalization source begin ---
% archon:physics
% archon:covers IPhO2026Problems/problem_IPhO_2026_2_C_1.lean
% archon:source-report reports/ipho_2026/problem_IPhO_2026_2_C_1.source.json
% archon:problem-id IPhO_2026_2
% archon:part-id C.1

\chapter{Physics problem IPhO\_2026\_2\_C\_1}
\label{ch:IPhO2026Problems_problem_IPhO_2026_2_C_1}

\paragraph{Problem source.}
For the half-cylindrical mirror of radius R, ray A is incident at angle theta
and its reflected line is y = m\_A*x + b\_A.  A neighboring parallel ray B is
incident at theta + Delta theta, with Delta theta much smaller than theta, and
its reflected line is y = m\_B*x + b\_B.  The envelope/intersection of neighboring
rays forms the caustic.  Use Figure 2g and its coordinate convention.

Current subquestion:
Write the slope m\_A and intercept b\_A of reflected ray A in terms of theta and R.

\paragraph{Current subquestion.}
Write the slope m\_A and intercept b\_A of reflected ray A in terms of theta and R.

\paragraph{Recorded answer/context.}
m\_A = cot(2*theta), and b\_A = R/(2*cos(theta)).

\paragraph{Figure/image path.}
/root/proposal\_for\_physic/science-mango/ipho\_2026\_source/image/T2\_page-4.png

\paragraph{Formalization target.}
create a compiling Lean file with sorry bodies at `IPhO2026Problems/problem\_IPhO\_2026\_2\_C\_1.lean`. The Lean declarations must preserve the physical quantities, dimensions or dimensional roles, figure labels, governing-law hypotheses, and final relation expressed by this problem.
Use Mathlib/Physlib names found through LeanExplore where available. If a physics API is missing, introduce faithful local abstractions rather than scalar placeholder aliases.

\begin{definition}[PlanePoint]
  \label{decl:physics:IPhO_2026_2_C_1:PlanePoint}
  \lean{IPhO2026Problems.IPhO2026_2_C_1.PlanePoint}
  A point in the Cartesian coordinate convention of Figure 2g.
\end{definition}
\begin{proof}
  This is the defining declaration; unfolding it gives the stated typed data or relation.
\end{proof}

\begin{definition}[HalfCylindricalMirror]
  \label{decl:physics:IPhO_2026_2_C_1:HalfCylindricalMirror}
  \lean{IPhO2026Problems.IPhO2026_2_C_1.HalfCylindricalMirror}
  The half-cylindrical mirror in Figure 2g, represented by the radius of its upper semicircular cross-section.
\end{definition}
\begin{proof}
  This is the defining declaration; unfolding it gives the stated typed data or relation.
\end{proof}

\begin{definition}[SlopeInterceptRay]
  \label{decl:physics:IPhO_2026_2_C_1:SlopeInterceptRay}
  \lean{IPhO2026Problems.IPhO2026_2_C_1.SlopeInterceptRay}
  A reflected ray represented in Figure 2g by “y = slope * x + intercept”.
\end{definition}
\begin{proof}
  This is the defining declaration; unfolding it gives the stated typed data or relation.
\end{proof}

\begin{definition}[OnUpperHalfMirror]
  \label{decl:physics:IPhO_2026_2_C_1:OnUpperHalfMirror}
  \lean{IPhO2026Problems.IPhO2026_2_C_1.OnUpperHalfMirror}
  \uses{decl:physics:IPhO_2026_2_C_1:PlanePoint, decl:physics:IPhO_2026_2_C_1:HalfCylindricalMirror}
  A point lies on the upper semicircular mirror centered at the origin.
\end{definition}
\begin{proof}
  This is the defining declaration; unfolding it gives the stated typed data or relation.
\end{proof}

\begin{definition}[LiesOnRayLine]
  \label{decl:physics:IPhO_2026_2_C_1:LiesOnRayLine}
  \lean{IPhO2026Problems.IPhO2026_2_C_1.LiesOnRayLine}
  \uses{decl:physics:IPhO_2026_2_C_1:PlanePoint, decl:physics:IPhO_2026_2_C_1:SlopeInterceptRay}
  Incidence of a point on the slope-intercept line supporting a ray.
\end{definition}
\begin{proof}
  This is the defining declaration; unfolding it gives the stated typed data or relation.
\end{proof}

\begin{definition}[ObeysSpecularReflection]
  \label{decl:physics:IPhO_2026_2_C_1:ObeysSpecularReflection}
  \lean{IPhO2026Problems.IPhO2026_2_C_1.ObeysSpecularReflection}
  The equal-angle law of specular reflection, written in terms of oriented direction angles and the tangent line at the impact point.
\end{definition}
\begin{proof}
  This is the defining declaration; unfolding it gives the stated typed data or relation.
\end{proof}

\begin{theorem}[Physics formalization target]
\label{thm:physics:IPhO_2026_2_C_1:target}
\lean{IPhO2026Problems.IPhO2026_2_C_1.rayA_slope_and_intercept}
\uses{decl:physics:IPhO_2026_2_C_1:PlanePoint, decl:physics:IPhO_2026_2_C_1:HalfCylindricalMirror, decl:physics:IPhO_2026_2_C_1:SlopeInterceptRay, decl:physics:IPhO_2026_2_C_1:OnUpperHalfMirror, decl:physics:IPhO_2026_2_C_1:LiesOnRayLine, decl:physics:IPhO_2026_2_C_1:ObeysSpecularReflection}
The assigned autoformalize agent should translate this physics problem into Lean declarations in the covered file, with theorem and lemma proof bodies written as `by sorry`.
\end{theorem}
\begin{proof}
This is an autoformalization task, not a proof task. Produce statements that can later be proved by the physics prover without weakening the physical model.
\end{proof}
% --- Archon physics formalization source end ---
